\label{chap:analysis}
\subsection{Discussion of Experimental Results}

This chapter provides a critical interpretation of the experimental findings presented in Chapter~5, examining the results from the perspective of the five research objectives, the identified research gaps, and the practical implications for industrial deployment. The discussion synthesizes the quantitative evidence from four experimental components --- T-GCN forecasting, supervised classification and regression, non-heuristic agentic pipeline, and heuristic-augmented pipeline --- and relates the findings to the broader literature on predictive maintenance.

\subsection{Spatiotemporal Forecasting: T-GCN vs.\ LSTM Baselines}

The T-GCN with cross-correlation adjacency achieves the best average R\textsuperscript{2} of 0.7702 across five sensor channels, outperforming both the OptimizedLSTM (R\textsuperscript{2} = 0.6853, $\Delta = +12.4\%$) and the Hybrid LSTM-GCN (R\textsuperscript{2} = 0.7134, $\Delta = +8.0\%$). This result confirms the core hypothesis underlying the T-GCN design: that modeling the spatial dependencies between machines through a cross-correlation lag graph provides information that temporal-only recurrent architectures cannot recover from the time series alone.

The improvement is most pronounced for humidity prediction (R\textsuperscript{2}: 0.9080 vs.\ 0.8761 for LSTM, $\Delta = +3.6\%$) and energy consumption (R\textsuperscript{2}: 0.6581 vs.\ 0.5522, $\Delta = +19.2\%$). The energy channel improvement is particularly notable because energy consumption exhibits near-zero linear correlation with other sensors individually, yet benefits substantially from graph-based cross-machine information propagation. This suggests that energy patterns across machines carry latent correlated structure --- possibly reflecting shared electrical load or production scheduling dependencies --- that the cross-correlation graph successfully captures.

The pressure sensor channel is the most difficult to forecast across all three models (R\textsuperscript{2} range: 0.4254--0.5002). Pressure has only 401 unique values in the dataset, suggesting quantized setpoint-driven dynamics that follow discrete state transitions rather than smooth continuous degradation trajectories. These step-function dynamics are inherently harder to model with gradient-based temporal architectures trained on mean-squared error, and future work should investigate specialized architectures for semi-discrete sensor channels.

The Hybrid LSTM-GCN achieves intermediate performance (R\textsuperscript{2} = 0.7134) with the shortest training time (47.89 seconds), offering a practical trade-off when computational resources are limited. However, its advantage over the pure LSTM is modest ($\Delta = +4.1\%$), suggesting that the naive concatenation of LSTM and GCN outputs without the structured lag-based graph topology of the T-GCN fails to fully realize the benefit of graph-based cross-machine reasoning. The T-GCN's principled construction of directed, lag-weighted edges from temporal cross-correlations is the key distinguishing design choice that delivers the superior performance.

Training time for the T-GCN (111.28 seconds) is approximately twice that of the LSTM (52.56 seconds), attributable to the graph convolution operations over 538 edges at each forward pass. In an operational setting where models are retrained periodically rather than continuously, this overhead is acceptable given the performance benefit. The \texttt{LearningLoop}'s retraining trigger, which requires a minimum of 5 feedback records before initiating retraining, provides a natural mechanism for controlling retraining frequency in proportion to the rate of observed distributional shift.

\subsection{Supervised Classification and Regression: Ensemble Models and TabNet}

The supervised learning experiment reveals a clear stratification of task difficulty across the six prediction targets. The anomaly flag and downtime risk binary classification targets are near-perfectly discriminable by all tested models, with LightGBM achieving F1 = 0.9990 and AUC-ROC = 0.9999. These results warrant careful interpretation: near-perfect classification performance suggests that these labels may be deterministically derivable from other features in the dataset --- specifically, the strong correlation ($r \approx 1.0$) between anomaly flag and downtime risk implies that they are effectively the same binary signal. In a real-world deployment context, a model that near-perfectly predicts anomaly flag or downtime risk from the same sensor reading may be learning a data leakage relationship rather than a genuinely predictive one. Future work should verify the temporal independence of these labels relative to the input features.

Machine status classification (best F1 = 0.712, XGBoost) and maintenance required classification (best F1 = 0.872, XGBoost) present more meaningful learning challenges that are representative of real-world difficulty. The failure type five-class classification (best F1 = 0.880, XGBoost) is the most challenging classification task, reflecting the severe class imbalance between Normal operation (91.90\%) and minority failure types (1.01--3.13\%). XGBoost's consistent superiority for the most challenging classification targets (machine status, failure type, maintenance required) is consistent with the broader literature demonstrating gradient-boosted tree advantages on tabular data with class imbalance and mixed feature types.

TabNet consistently underperforms relative to gradient-boosted ensemble methods across all targets, consistent with recent empirical comparisons that show TabNet's attention-based selection mechanism provides less benefit on tabular datasets with moderate feature counts and clear feature importances. The 50-machine dataset, with only five raw sensor channels and a modest set of engineered features, may not provide sufficient feature complexity to benefit from the sequential attention architecture that distinguishes TabNet.

The RUL regression results (best R\textsuperscript{2} = 0.189, CatBoost) deserve particular scrutiny. An R\textsuperscript{2} of 0.189 indicates that the five sensor channels and derived features in a single-timestep snapshot explain only 18.9\% of the variance in the RUL label. This is expected: the near-uniform RUL distribution (mean 234.27 hours, standard deviation 150.06 hours) means that a reading with, say, temperature = 75\textdegree C and vibration = 50 could correspond to any RUL value between 1 and 499 hours depending on the machine's degradation history and initial condition --- information that is absent from a single-snapshot feature vector. The T-GCN forecasting architecture addresses this fundamental limitation by modeling the temporal trajectory of sensor evolution, which carries the degradation history information that static snapshots lack. This result reinforces the methodological rationale for the T-GCN component in the QASAMAP pipeline.

\subsection{Agentic AI Pipeline Without Quantum Enhancement}

The non-quantum agentic pipeline processes 11 events across 7 machines and produces active maintenance decisions for only 1 out of 7 machines (Machine 32, classical risk = 0.300 MEDIUM), while the remaining 6 machines receive passive monitoring via the \texttt{auto-LOW-skip} rule. This 14.3\% machine coverage reflects a fundamental limitation of classical threshold-based risk assessment: machines operating below the MEDIUM risk boundary are invisible to the active decision-making pipeline regardless of whether their sensor patterns carry subtle multi-sensor risk combinations detectable by the heuristic layer.

The Monitoring Agent's assessment for Machine 32 --- bearing wear with vibration and humidity WARN flags at 48-hour urgency --- is diagnostically reasonable given the classical evidence (vibration = 67.09, WARN threshold = 55; humidity = 74.35, WARN threshold = 68). The Planning Agent's recommendation of enhanced monitoring at P4 rather than immediate inspection reflects the conservative classical risk assessment (risk = 0.300, boundary of MEDIUM level) and the high RUL (446 hours), which together suggest that urgent intervention is not yet warranted from a purely classical perspective. This conservative behavior is appropriate for a classical system but represents a potential gap when sub-threshold heuristic risk signals exist for other machines.

The limited feedback from the non-quantum run (1 feedback event vs.\ the minimum of 5 required for retraining) reflects the low active-decision rate of the pipeline. This is an important operational consideration: a classical-only agentic system that routes the majority of events to passive monitoring will accumulate insufficient feedback for learning loop adaptation, limiting the system's ability to improve through experience. The heuristic-augmented pipeline's higher active-decision rate (10 active decisions out of 11 events) generates richer feedback that supports more effective learning loop operation.

\subsection{Heuristic-Augmented Agentic Pipeline}

The quantum-inspired classical heuristic layer changes the decision profile of the agentic pipeline in the 7-machine demonstration. The QUBO local search identifies Machine~32 as the highest-priority machine (heuristic risk = 0.5144, energy = $-$0.5144), driven by its humidity Pauli-Z value of $-$0.8704 --- the sensor value closest to the heuristic danger threshold of $-1$ in the entire experimental batch. This diagnosis aligns with the classical evidence (humidity WARN at 74.35\%) but provides a more precise, quantitative danger signal: the Pauli-Z value of $-$0.8704 places Machine~32's humidity sensor at 93.5\% of the way from the safe zone ($\langle Z \rangle = +1$) to the danger zone ($\langle Z \rangle = -1$), a level of precision that classical binary WARN/CRITICAL threshold crossing does not provide. \textbf{Algorithmic disclosure:} These values are computed by the classical \texttt{energy\_distance\_anomaly\_score} function; no quantum circuit is executed.

The greedy graph-message-passing scheduler partitions machines such that machines with higher heuristic risk scores and QUBO selections are prioritized for Slot~A. The heuristic override mechanism triggers for 5 out of 7 machines (M-1, M-2, M-3, M-4, M-32, M-40), escalating them from passive monitoring to active inspection or scheduled maintenance despite classical LOW or borderline MEDIUM risk levels. Machine~47, classified as QUANTUM-MARGINAL ($\langle Z \rangle_c = +0.3608$, heuristic risk = 0.3196) and ranked last by QUBO, correctly receives passive monitoring --- demonstrating that the heuristic layer does not indiscriminately escalate all machines, but rather applies a principled risk quantification that preserves the selectivity essential for operational credibility.

The dramatic improvement in machine coverage (from 14.3\% to 85.7\%) is the most striking difference in the 7-machine demonstration. This expansion in active decision coverage, achieved without changing the classical risk thresholds, demonstrates that the heuristic layer functions as a sensitivity amplification mechanism that detects genuine risk signals at sensor value combinations below the classical warning boundaries. \textbf{Critical caveat:} This demonstration involves a small, fixed sample of 7 machines. On the held-out 25-machine test split, the agentic pipeline --- with or without heuristic enhancement --- achieves F1=0.00, identical to the rule-based baseline. The heuristic override improves demonstrative coverage but does not solve the fundamental generalisation problem on unseen data.

The 72.7\% false positive rate recorded in the heuristic-augmented feedback log (compared to 0\% for the non-augmented run) requires careful interpretation. In the context of this experiment, a false positive is a heuristic-override-driven inspection action where no classical-level failure was confirmed in the immediate feedback. This elevated false positive rate reflects the heuristic layer's inherently conservative design: it is calibrated to minimize false negatives (missed failures) at the cost of some false positives (unnecessary inspections). The Learning Loop addresses this explicitly through its threshold relaxation mechanism, which is triggered when the false positive rate exceeds 40\%. In the experimental run, a false positive rate of 72.7\% would immediately trigger threshold relaxation in the next decision cycle, beginning the adaptive process of calibrating the heuristic override sensitivity to the specific operating context of the monitored machine pool.

\subsection{Post-Audit Comparative Analysis: Interpretation and Significance}

The post-audit comparative analysis replaces the pre-audit KPI tables, which were found to contain fabricated metrics derived from label-leaking experimental procedures. The corrected evaluation reports two scopes: a held-out test split (25 machines) and a full-dataset comparison (50 machines).

On the held-out test split, both the Rule-Based and AgenticAI pipelines achieve F1=0.00, precision=0.00, and recall=0.00. This result is sobering: despite the rich capabilities of the nine-agent pipeline (root cause diagnosis, explainability, correlation, consensus, scheduling, pattern detection), its conservative static thresholds completely fail to detect anomalies in the held-out test set. The Fleet-Relative detector (F1=0.30) and Sliding-Window detector (F1=0.48) establish stronger baselines by exploiting machine-relative deviations and temporal context, respectively, suggesting that future agentic pipelines should integrate these detection strategies rather than relying solely on static thresholds.

On the full 50-machine dataset (pre-audit run, where ground-truth fields were previously leaked), the AgenticAI pipeline achieves perfect precision (1.00) but very low recall (0.20), yielding F1=0.33. This indicates that the pipeline is highly specific --- when it flags a machine, it is almost always correct --- but it misses 80\% of actual failures. By contrast, the neural baselines (Transformer, LSTM, GNN, RL-QLearn, FedLearn, TransferL) achieve F1 scores of 0.73--0.76 but produce \textbf{zero true negatives} (TN=0 or TN=1), revealing systematic over-prediction: they classify nearly every machine as anomalous. This over-prediction would be operationally catastrophic in a real factory, generating maintenance orders for the entire fleet. The AgenticAI pipeline's high specificity is therefore a genuine operational advantage, but its low sensitivity is a critical deployment blocker.

The latency comparison exposes a fundamental trade-off. The AgenticAI pipeline requires 5,771.8~seconds ($\approx$1.6~hours) per prediction cycle, versus milliseconds for neural baselines. This latency is driven by nine sequential LLM API calls per machine event. The LLM benchmark (Table~\ref{tab:llm_benchmark}) reveals that even the fastest model (DeepSeek-V4, 20.8~s average) cannot reduce the total cycle below $\sim$200~seconds per machine, and the Reporting Agent consistently fails structured output (0\% JSON rate), forcing manual parsing of free-text alerts. These reliability and latency issues are discussed further in Chapter~\ref{chap:analysis}.

The quantum-inspired heuristic layer's contribution must be interpreted with full disclosure. The \texttt{PauliZRiskEncoder}, \texttt{QuantumAnnealingOptimizer}, and \texttt{QAOAScheduler} are classical heuristics implemented in \texttt{classical\_optimization.py}. They preserve the mathematical structure of their quantum counterparts (QUBO/Ising, Pauli-Z expectation, QAOA Hamiltonian) for future hardware migration, but the current experimental results reflect classical CPU execution. Claims of ``quantum'' performance advantages in earlier thesis drafts have been removed in the post-audit rewrite (Ver13, 2026-05).

\subsection{Implications for Industrial Deployment}

The experimental results collectively support three practical implications for industrial deployment of the QASAMAP framework. First, the T-GCN's superior forecasting performance over LSTM baselines, combined with the \texttt{AnchorBufferManager}'s ability to handle partial sensor instrumentation, makes it feasible to deploy T-GCN-based streaming inference in facilities where not all machines are fully instrumented. The imputation of missing sensor values from T-GCN predictions reduces the minimum sensor infrastructure required for high-quality streaming inference.

Second, the near-perfect performance of tree-based ensemble models on anomaly flag and downtime risk classification, combined with the meaningful performance on failure type and maintenance required classification, suggests that a hybrid decision architecture --- using fast ensemble model inference for initial screening and routing --- could substantially reduce the LLM agent invocation rate, improving system throughput and reducing operational cost. Specifically, events classified as anomaly flag = 1 with high confidence by LightGBM could be directly escalated without requiring the full nine-agent reasoning chain.

Third, the heuristic override mechanism's high false positive rate in the current experimental context, while addressed by the Learning Loop, highlights the importance of the adaptive threshold calibration for operational credibility. An uncalibrated heuristic system that generates inspections for 85.7\% of machines simultaneously would overwhelm maintenance resources and erode operator trust. The Learning Loop's threshold relaxation mechanism is therefore not merely a performance optimization but an operational necessity for production deployment. Future deployments should include a calibration period of sufficient duration to accumulate the feedback needed for effective threshold adaptation before moving to full autonomous operation.